% ============================================================
% PAGE 7 — Module 3: Detection Head & Optimized Loss Functions
% ============================================================

\subsection{Module 3: Detection Head and Optimized Loss
Functions}

The detection head is the final computational stage of
FusionDet: it consumes the four-level feature pyramid
$\{\hat{\mathbf{N}}_l\}_{l=2}^{5}$ produced by the neck and
emits, for every spatial position at every scale, a joint
prediction of object presence, class identity, and bounding-box
geometry. FusionDet adopts an \textit{anchor-free, decoupled}
head design in which classification and regression are handled
by independent sub-networks---an architectural choice
empirically shown to eliminate the task-interference that arises
when a single shared feature branch must simultaneously
represent discriminative appearance features (for
classification) and precise geometric offsets (for regression)
\cite{ge2021yolox}.

\subsubsection{Decoupled Head Structure}

For each pyramid level $l$, the head receives the
attention-refined feature map
$\hat{\mathbf{N}}_l \in \mathbb{R}^{C \times H_l \times W_l}$
and passes it through two parallel branches:

\begin{itemize}
  \item \textbf{Classification branch:} two $3\times3$
    convolutional layers with SiLU activations, followed by a
    $1\times1$ projection to $C_{\text{cls}}$ channels, yielding
    per-location class probability logits.
  \item \textbf{Regression branch:} two $3\times3$
    convolutional layers with SiLU activations, followed by a
    $1\times1$ projection to four channels encoding the
    center-offset and log-scale box parameters
    $(\Delta x, \Delta y, \Delta w, \Delta h)$.
\end{itemize}

An additional lightweight \textbf{objectness branch} (one
$1\times1$ layer) produces a binary foreground score for each
location, which is used during inference to suppress predictions
at background locations before applying non-maximum
suppression. The decoupled architecture adds only ${\sim}0.1$~M
parameters per pyramid level relative to a coupled head, while
consistently improving mAP by $1.0$--$1.5$~pp in our
experiments.

\subsubsection{Anchor-Free Target Assignment}

Unlike anchor-based heads that pre-define a discrete set of
aspect ratios and scales at each grid cell, FusionDet directly
predicts box offsets from every foreground grid cell. During
training, a ground-truth box $b = (x, y, w, h)$ is assigned to
level $l$ if $\max(w, h)$ falls within a scale range
$[s_l, s_{l+1})$ pre-defined per level, and to spatial cell
$(i, j)$ if the cell center lies inside the box. This
SimOTA-inspired \cite{ge2021yolox} dynamic assignment resolves
ambiguity for small or heavily overlapping boxes by selecting
the top-$k$ candidate cells per ground-truth based on joint
classification and regression cost.

\subsubsection{Composite Training Loss}

The full training objective for FusionDet is a weighted sum of
three terms:
\begin{equation}
  L_{\text{total}} =
    \lambda_{\text{cls}}\,L_{\text{cls}}
    + \lambda_{\text{box}}\,L_{\text{CIoU}}
    + \lambda_{\text{obj}}\,L_{\text{obj}}
  \label{eq:loss_total}
\end{equation}
with $\lambda_{\text{cls}} = 1.0$, $\lambda_{\text{box}} = 5.0$,
and $\lambda_{\text{obj}} = 1.0$ (validated by grid search).

\paragraph{Classification Loss.}
We adopt Varifocal Loss (VFL), which generalizes the standard
binary cross-entropy by assigning soft quality-aware labels and
weighting the loss non-uniformly across foreground and
background samples. For a predicted class-score vector
$\hat{\mathbf{p}}$ and a quality-weighted target
$\mathbf{q}$ (set to the predicted IoU for positives and zero
for negatives), VFL is:
\begin{equation}
  L_{\text{cls}} = -\sum_{i}
  \Bigl[
    q_i \log \hat{p}_i
    + \alpha\,(1 - q_i)\,\hat{p}_i^{\gamma}
      \log(1 - \hat{p}_i)
  \Bigr]
  \label{eq:vfl}
\end{equation}
where $\alpha = 0.75$ and $\gamma = 2.0$ are focal parameters
that down-weight the contribution of easy negatives.

\paragraph{Bounding-Box Regression Loss (CIoU).}
Intersection-over-Union (IoU) is the canonical measure of box
overlap. Simple $\ell_1$ or $\ell_2$ regression losses are not
scale-invariant and do not directly optimize IoU. Complete IoU
Loss \cite{zheng2020ciou} addresses both issues by
decomposing the regression penalty into three geometrically
meaningful terms.

Let $B = (x, y, w, h)$ denote the predicted box and
$B^* = (x^*, y^*, w^*, h^*)$ the matched ground-truth box.
Define $\rho^2$ as the squared Euclidean distance between
their centers:
\begin{equation}
  \rho^{2} =
  (x - x^{*})^{2} + (y - y^{*})^{2}
  \label{eq:ciou_dist}
\end{equation}
and let $c$ denote the diagonal length of the smallest
enclosing box of $B$ and $B^*$. The aspect-ratio consistency
term is:
\begin{equation}
  v = \frac{4}{\pi^{2}}
  \left(
    \arctan\frac{w^{*}}{h^{*}}
    - \arctan\frac{w}{h}
  \right)^{2}
  \label{eq:ciou_v}
\end{equation}
with the trade-off coefficient:
\begin{equation}
  \alpha_v =
  \frac{v}{(1 - \text{IoU}) + v}
  \label{eq:ciou_alpha}
\end{equation}
The CIoU loss is then:
\begin{equation}
  L_{\text{CIoU}} =
  1 - \text{IoU}
  + \frac{\rho^{2}}{c^{2}}
  + \alpha_v\,v
  \label{eq:ciou_loss}
\end{equation}
The first term penalizes insufficient overlap, the second
penalizes center-point misalignment relative to the enclosing
diagonal, and the third penalizes aspect-ratio inconsistency.
Together they provide a comprehensive, differentiable proxy for
the localization quality of each predicted box
\cite{zheng2020ciou,rezatofighi2019giou}.

\paragraph{Objectness Loss.}
The objectness branch is trained with binary cross-entropy,
where the positive label for cell $(i,j)$ at level $l$ is set
to the IoU between the predicted box and the assigned
ground-truth box (capped at 1.0), encoding the notion that a
cell's objectness score should reflect \textit{how well} the
box it predicts covers the ground truth rather than merely
\textit{whether} any object is present:
\begin{equation}
  L_{\text{obj}} = -\sum_{i,j,l}
  \Bigl[
    t_{ij}^{l} \log o_{ij}^{l}
    + (1 - t_{ij}^{l})\log(1 - o_{ij}^{l})
  \Bigr]
  \label{eq:obj_loss}
\end{equation}
where $t_{ij}^{l} = \text{IoU}(\hat{B}_{ij}^l, B^*)$ for
positive cells and $t_{ij}^{l} = 0$ for negatives.

\subsubsection{Summary}

Table~\ref{tab:head_summary} summarizes the parameters and
computational cost of the decoupled head at each pyramid level
for an input resolution of $640\times640$.

\begin{table}[h]
\centering
\caption{Detection Head Cost per Pyramid Level ($C=256$,
         Input $640{\times}640$).}
\label{tab:head_summary}
\renewcommand{\arraystretch}{1.1}
\begin{tabular}{@{}lcccc@{}}
\toprule
Level & Stride & Spatial & Params (M) & GFLOPs \\
\midrule
$\hat{\mathbf{N}}_2$ & 4  & $160{\times}160$ & 0.66 & 1.08 \\
$\hat{\mathbf{N}}_3$ & 8  & $80{\times}80$   & 0.66 & 0.27 \\
$\hat{\mathbf{N}}_4$ & 16 & $40{\times}40$   & 0.66 & 0.07 \\
$\hat{\mathbf{N}}_5$ & 32 & $20{\times}20$   & 0.66 & 0.02 \\
\midrule
\textbf{Total}       &    &                  & \textbf{2.64} & \textbf{1.44} \\
\bottomrule
\end{tabular}
\end{table}

The decoupled head's total cost of $1.44$~GFLOPs represents
less than $3\%$ of FusionDet's overall inference budget, yet
its decoupled design and quality-aware loss functions are
responsible for the majority of the accuracy advantage over
coupled-head baselines, as confirmed by the ablation study
in Section~\ref{sec:results}.
